\section{Governance in Practice}

\begin{frame}{Organisational Practice}

    \begin{columns}[T]
        \begin{column}{0.49\textwidth}
            \textbf{Review before release}
            \begin{itemize}\small
                \setlength\itemsep{0.3em}
                \item A named reviewer who is \textbf{not} the model's author signs off the card
                      and the disaggregated results.
                \item Scale the review to the stakes: a recommender for internal search does not
                      need a committee; a credit model does.
                \item Record the decision \emph{and the trade-off accepted}, so it can be revisited.
            \end{itemize}

            \vspace{0.7em}
            \textbf{Red-teaming}
            \begin{itemize}\small
                \setlength\itemsep{0.3em}
                \item Deliberately try to make the system fail for a specific group or input type,
                      before users do.
                \item Give the red team the failure taxonomy from this lecture as a starting
                      checklist.
            \end{itemize}
        \end{column}
        \begin{column}{0.49\textwidth}
            \textbf{After release}
            \begin{itemize}\small
                \setlength\itemsep{0.3em}
                \item \textbf{Incident reporting.} A defined path for ``the model did something
                      harmful,'' with an owner and a deadline --- not a support ticket queue.
                \item \textbf{Contestability.} A person affected by a decision can ask for it to be
                      reviewed by a human, and the request is logged.
                \item \textbf{Scheduled re-audit.} Put the date in the model card and treat it like
                      a certificate expiry.
                \item \textbf{A kill switch.} Know in advance how you turn the model off and what
                      the fallback is.
            \end{itemize}
        \end{column}
    \end{columns}

    \vspace{0.8em}
    \begin{block}{}
        \centering
        \small The common failure is not the absence of a policy. It is a policy with no owner and
        no date.
    \end{block}
\end{frame}

